% =====================================================================
%  Capítulo 14 — Aplicaciones móviles de seguridad
% =====================================================================
\chapter{Aplicaciones móviles de seguridad}
\citacap{La tecnología debe ser un aliado de tu atención, nunca un sustituto de ella.}{Anónimo}

El teléfono móvil es hoy una herramienta de autoprotección de primer orden: permite pedir ayuda, compartir ubicación en tiempo real y documentar incidentes. Estas son las aplicaciones recomendadas en España (todas gratuitas y oficiales):

\section{AlertCops (Ministerio del Interior) --- la imprescindible}

AlertCops es la aplicación oficial de seguridad ciudadana de la Secretaría de Estado de Seguridad (Ministerio del Interior), disponible en toda España. Proporciona un canal directo con Policía Nacional y Guardia Civil, con geolocalización de la alerta para que acuda la patrulla más próxima. Sus funciones principales:

\begin{itemize}
  \item \textbf{Alertas con chat:} comunicar cualquier hecho del que seas víctima o testigo, adjuntando fotos, vídeos, audios o tu localización, con respuesta inmediata.
  \item \textbf{Botón SOS:} con protección reforzada para colectivos vulnerables (entre ellos, víctimas de violencia de género). Pulsándolo repetidamente (cinco veces en menos de seis segundos) envía una alerta urgente al centro policial más cercano con tu posición y una grabación de audio de 10 segundos para que los agentes valoren la gravedad.
  \item \textbf{Guardián:} comparte tu ubicación en tiempo real con contactos de confianza o con las Fuerzas de Seguridad.
  \item \textbf{Acompáñame:} para trayectos en los que quieras sentirte acompañada (volver a casa de noche, una cita con un desconocido).
  \item \textbf{Llamada simulada:} genera una llamada falsa que te da excusa para salir de una situación incómoda.
  \item \textbf{Traductor automático} para comunicarse con la policía en más de cien idiomas, y alertas de seguridad y emergencia de tu zona.
\end{itemize}

Descarga oficial (comprueba que el desarrollador es \enquote{Secretaría de Estado de Seguridad - Ministerio del Interior}):

\begin{itemize}
  \item iOS (App Store): \url{https://apps.apple.com/es/app/alertcops/id1273718252}
  \item Android (Google Play): \url{https://play.google.com/store/apps/details?id=com.alertcops4.app}
  \item Web oficial: \url{https://alertcops.ses.mir.es}
\end{itemize}

Tras descargarla, regístrate con tu nombre y número de móvil (validación por SMS) y configura el Botón SOS y tus contactos guardianes ANTES de necesitarlos.

\section{Otras aplicaciones y servicios útiles}

\begin{itemize}
  \item \textbf{My112 / apps oficiales 112 autonómicas:} al llamar a emergencias envían tu posición exacta al centro de coordinación, crucial si no sabes dónde estás. Busca la app 112 oficial de tu comunidad autónoma en las tiendas oficiales.
  \item \textbf{016 Online:} además del teléfono 016, el servicio estatal contra la violencia de género atiende por WhatsApp (\recurso{600 000 016}) y correo (\href{mailto:016-online@igualdad.gob.es}{016-online@igualdad.gob.es}), útil cuando hablar en voz alta es peligroso.
  \item \textbf{Compartir ubicación del propio sistema:} tanto iOS (Buscar/Find My) como Android (Google Maps -- Compartir ubicación) permiten que personas de confianza sigan tu trayecto en tiempo real sin instalar nada más.
  \item \textbf{Función de emergencia del teléfono:} aprende el gesto SOS de tu modelo (en muchos Android e iPhone, pulsar 5 veces el botón lateral llama a emergencias y avisa a tus contactos). Configúralo hoy.
\end{itemize}

\begin{consejo}
\textbf{Consejo:} ninguna app sustituye a la atención al entorno. El móvil en la mano mirando la pantalla te pone en \enquote{código blanco}; el móvil configurado en el bolsillo, con SOS y Guardián listos, es un arma defensiva.
\end{consejo}
